\documentclass{article}
\usepackage{amsmath,amssymb,amsthm,algorithm,algpseudocode}
\usepackage{fullpage}
\usepackage{hyperref}
\newtheorem{theorem}{Theorem}
\newtheorem{lemma}{Lemma}
\newtheorem{assumption}{Assumption}
\begin{document}

\section*{Quantized Stochastic Gradient Descent (QSGD)}

\section*{Mathematical Formulation}
Let $f:\mathbb{R}^d\rightarrow\mathbb{R}$ be a convex and $L$–smooth function.  
At iteration $t$, a stochastic gradient $g_t$ is drawn such that 
\[
\mathbb{E}[\,g_t\mid w_t\,]=\nabla f(w_t),\qquad 
\mathbb{E}[\|g_t\|^2\mid w_t]\le G^2 .
\]
A (possibly random) quantization operator $Q:\mathbb{R}^d\rightarrow\mathbb{R}^d$ satisfies  

\[
\boxed{\mathbb{E}[\,Q(g)\mid g]=g}\qquad\text{(unbiasedness)} \tag{1}
\]
and a variance bound  

\[
\boxed{\mathbb{E}[\|Q(g)-g\|^2\mid g]\le \omega\,\|g\|^2}\qquad\text{(relative variance)} \tag{2}
\]
for some constant $\omega\ge0$ (e.g. $\omega =\frac{d}{s}-1$ for $s$‑bit stochastic rounding).

The update rule of QSGD is
\[
\boxed{w_{t+1}=w_t-\eta_t\,Q(g_t)}\tag{3}
\]
where $\{\eta_t\}_{t\ge0}$ is a stepsize schedule (constant or diminishing).

\section*{Pseudocode}
\begin{algorithm}[H]
\caption{Quantized Stochastic Gradient Descent}
\begin{algorithmic}[1]
\State \textbf{Input:} initial point $w_0\in\mathbb{R}^d$, stepsizes $\{\eta_t\}$, quantizer $Q(\cdot)$
\For{$t=0,1,2,\dots,T-1$}
    \State Sample a stochastic gradient $g_t$ at $w_t$ (e.g. $g_t=\nabla f(w_t;z_t)$)
    \State Compute quantized gradient $\tilde g_t = Q(g_t)$
    \State $w_{t+1}\gets w_t-\eta_t \tilde g_t$
\EndFor
\State \textbf{Output:} $\displaystyle \bar w_T=\frac{1}{T}\sum_{t=0}^{T-1} w_t$ (optional averaging)
\end{algorithmic}
\end{algorithm}

\section*{Dry Run Example}
Consider $f(w)=(w-3)^2$ (so $d=1$), $w_0=0$, constant stepsize $\eta=0.1$, and a simple 2‑level unbiased quantizer:
\[
Q(g)=\begin{cases}
+g & \text{with prob. }1/2,\\[2mm]
-g & \text{with prob. }1/2 .
\end{cases}
\]
Thus $\mathbb{E}[Q(g)]=0$ when $g=0$ and $\mathbb{E}[Q(g)]=g$ for any $g$ (unbiased) and $\operatorname{Var}(Q(g))=g^2$ ($\omega=1$).

\begin{center}
\begin{tabular}{c|c|c|c}
$t$ & $g_t=\nabla f(w_t)$ & $\tilde g_t=Q(g_t)$ (sample) & $w_{t+1}=w_t-\eta\tilde g_t$ \\ \hline
0 & $2(w_0-3)=-6$ & suppose $+6$ & $w_1=0-0.1\cdot6=-0.6$ \\
1 & $2(w_1-3)=2(-3.6)=-7.2$ & suppose $-7.2$ & $w_2=-0.6-0.1(-7.2)=0.12$ \\
2 & $2(w_2-3)=2(-2.88)=-5.76$ & suppose $+5.76$ & $w_3=0.12-0.1\cdot5.76=-0.456$
\end{tabular}
\end{center}
Objective values: $f(w_1)=(-0.6-3)^2=12.96$, $f(w_2)= (0.12-3)^2=8.2944$, $f(w_3)=(-0.456-3)^2=11.938$.
The example shows stochastic fluctuation caused by quantization.

\newpage
\section*{Assumptions}
\begin{assumption}[Convexity and Smoothness]\label{ass:convex}
$f$ is convex and $L$–smooth, i.e.,
\[
f(y)\le f(x)+\langle\nabla f(x),y-x\rangle+\frac{L}{2}\|y-x\|^2,\qquad\forall x,y\in\mathbb{R}^d .
\]
\end{assumption}
\begin{assumption}[Unbiased Stochastic Gradient]\label{ass:sgd}
\[
\mathbb{E}[\,g_t\mid w_t]=\nabla f(w_t),\qquad 
\mathbb{E}[\|g_t\|^2\mid w_t]\le G^2 .
\]
\end{assumption}
\begin{assumption}[Unbiased Quantization with Bounded Relative Variance]\label{ass:quant}
Quantizer $Q$ satisfies (1)–(2) with a constant $\omega\ge0$.
\end{assumption}
\begin{assumption}[Stepsize]\label{ass:stepsize}
We use a diminishing stepsize $\eta_t=\frac{c}{\sqrt{t+1}}$ with $c>0$ chosen later.
\end{assumption}

\section*{Main Convergence Theorem}
\begin{theorem}\label{thm:main}
Under Assumptions \ref{ass:convex}–\ref{ass:stepsize}, let $w^*\in\arg\min f$.  Define the averaged iterate 
\[
\bar w_T=\frac{1}{T}\sum_{t=0}^{T-1} w_t .
\]
Then
\[
\boxed{\mathbb{E}\bigl[f(\bar w_T)-f(w^*)\bigr]\le
\frac{\|w_0-w^*\|^2}{2c\sqrt{T}}+\frac{c(L+L\omega+1)G^2}{\sqrt{T}} } .
\]
Consequently $\mathbb{E}[f(\bar w_T)-f(w^*)]=\mathcal{O}\!\left(\tfrac{1}{\sqrt{T}}\right)$.
\end{theorem}

\section*{Theoretical Properties}
\begin{itemize}
\item \textbf{Stability.} The recursion (3) together with (2) guarantees that the variance of the update does not explode; the factor $(1+\omega)$ appears only as a multiplicative constant.
\item \textbf{Hyperparameter Sensitivity.} The optimal constant $c$ balances the bias term $\|w_0-w^*\|^2/(2c\sqrt{T})$ and the variance term $c(L+L\omega+1)G^2/\sqrt{T}$. Minimizing the bound yields $c^\star=\sqrt{\|w_0-w^*\|^2/(2(L+L\omega+1)G^2)}$.
\item \textbf{Comparison to Plain SGD.} Setting $\omega=0$ recovers the standard SGD bound (Theorem 1 in Example 1). Quantization incurs a penalty proportional to $\omega$, i.e. the coarser the quantization, the larger the constant.
\end{itemize}

\section*{Discussion}
The theorem shows that unbiased quantization preserves the optimal $\mathcal{O}(1/\sqrt{T})$ rate for convex smooth objectives. The proof follows the classical SGD analysis with an additional variance term stemming from (2). The bound is tight up to constants because the worst‑case variance of $Q(g)$ equals $(1+\omega)\|g\|^2$.

\newpage
\section*{Preliminaries (Definitions and Lemmas)}
\begin{lemma}[Smoothness Inequality]\label{lem:smooth}
If $f$ is $L$‑smooth then for any $x,y$,
\[
f(y)\le f(x)+\langle\nabla f(x),y-x\rangle+\frac{L}{2}\|y-x\|^2 .
\]
\end{lemma}
\begin{lemma}[Convexity Inequality]\label{lem:convex}
If $f$ is convex then for any $x,y$,
\[
f(x)-f(y)\le\langle\nabla f(x),x-y\rangle .
\]
\end{lemma}
\begin{lemma}[Variance Decomposition]\label{lem:var}
For any random vector $Z$ and any $\sigma$–algebra $\mathcal{F}$,
\[
\mathbb{E}\bigl[\|Z\|^2\mid\mathcal{F}\bigr]
= \|\mathbb{E}[Z\mid\mathcal{F}]\|^2
+ \mathbb{E}\bigl[\|Z-\mathbb{E}[Z\mid\mathcal{F}]\|^2\mid\mathcal{F}\bigr].
\]
\end{lemma}

\section*{Main Proof (Step‑by‑Step Derivation)}
We prove Theorem \ref{thm:main}.  Throughout we condition on the filtration
$\mathcal{F}_t:=\sigma(w_0,\dots,w_t,g_0,\dots,g_{t-1},\tilde g_0,\dots,\tilde g_{t-1})$.

\noindent\textbf{[Step 1]}  Start from the squared distance recursion induced by (3):
\[
\|w_{t+1}-w^*\|^2
 =\|w_t-w^*-\eta_t\tilde g_t\|^2
 =\|w_t-w^*\|^2-2\eta_t\langle\tilde g_t,w_t-w^*\rangle+\eta_t^2\|\tilde g_t\|^2 .
\tag{4}
\]

\noindent\textbf{[Step 2]}  Take conditional expectation w.r.t. $\mathcal{F}_t$ and use Lemma~\ref{lem:var} together with unbiasedness (1):
\[
\mathbb{E}\bigl[\langle\tilde g_t,w_t-w^*\rangle\mid\mathcal{F}_t\bigr]
 =\Big\langle\mathbb{E}[\tilde g_t\mid\mathcal{F}_t],\,w_t-w^*\Big\rangle
 =\langle g_t,w_t-w^*\rangle .
\tag{5}
\]

\noindent\textbf{[Step 3]}  For the second‑order term, apply Lemma~\ref{lem:var} again:
\[
\begin{aligned}
\mathbb{E}\bigl[\|\tilde g_t\|^2\mid\mathcal{F}_t\bigr]
&=\|\mathbb{E}[\tilde g_t\mid\mathcal{F}_t]\|^2
   +\mathbb{E}\bigl[\|\tilde g_t-\mathbb{E}[\tilde g_t\mid\mathcal{F}_t]\|^2\mid\mathcal{F}_t\bigr]\\
&=\|g_t\|^2+\mathbb{E}\bigl[\|Q(g_t)-g_t\|^2\mid g_t\bigr]
 \le (1+\omega)\|g_t\|^2 .
\end{aligned}
\tag{6}
\]

\noindent\textbf{[Step 4]}  Insert (5)–(6) into (4) and take full expectation:
\[
\mathbb{E}\|w_{t+1}-w^*\|^2
\le \mathbb{E}\|w_t-w^*\|^2
      -2\eta_t\mathbb{E}\langle g_t,w_t-w^*\rangle
      +\eta_t^2(1+\omega)\mathbb{E}\|g_t\|^2 .
\tag{7}
\]

\noindent\textbf{[Step 5]}  Use convexity Lemma~\ref{lem:convex} with $x=w_t$, $y=w^*$:
\[
\langle g_t,w_t-w^*\rangle
 =\langle \nabla f(w_t),w_t-w^*\rangle
 \ge f(w_t)-f(w^*) .
\tag{8}
\]

\noindent\textbf{[Step 6]}  Apply the bound $\mathbb{E}\|g_t\|^2\le G^2$ from Assumption~\ref{ass:sgd} and substitute (8) into (7):
\[
\mathbb{E}\|w_{t+1}-w^*\|^2
\le \mathbb{E}\|w_t-w^*\|^2
      -2\eta_t\mathbb{E}[f(w_t)-f(w^*)]
      +\eta_t^2(1+\omega)G^2 .
\tag{9}
\]

\noindent\textbf{[Step 7]}  Rearrange (9) to isolate the suboptimality term:
\[
2\eta_t\mathbb{E}[f(w_t)-f(w^*)]
\le \mathbb{E}\|w_t-w^*\|^2-\mathbb{E}\|w_{t+1}-w^*\|^2
      +\eta_t^2(1+\omega)G^2 .
\tag{10}
\]

\noindent\textbf{[Step 8]}  Sum (10) over $t=0,\dots,T-1$:
\[
2\sum_{t=0}^{T-1}\eta_t\mathbb{E}[f(w_t)-f(w^*)]
\le \|w_0-w^*\|^2
    + (1+\omega)G^2\sum_{t=0}^{T-1}\eta_t^2 .
\tag{11}
\]

\noindent\textbf{[Step 9]}  Use the stepsize choice $\eta_t = \dfrac{c}{\sqrt{t+1}}$.
Then
\[
\sum_{t=0}^{T-1}\eta_t = c\sum_{t=1}^{T}\frac{1}{\sqrt{t}}
    \ge c\int_{0}^{T}\frac{dt}{\sqrt{t+1}}
    = 2c\bigl(\sqrt{T+1}-1\bigr) \ge 2c\sqrt{T} -2c .
\tag{12}
\]
Similarly,
\[
\sum_{t=0}^{T-1}\eta_t^2 = c^2\sum_{t=1}^{T}\frac{1}{t}
    \le c^2\bigl(1+\ln T\bigr) \le 2c^2\ln T \quad\text{(for $T\ge e$)} .
\tag{13}
\]

\noindent\textbf{[Step 10]}  Divide (11) by $2\sum_{t=0}^{T-1}\eta_t$ and use the lower bound (12):
\[
\frac{1}{T}\sum_{t=0}^{T-1}\mathbb{E}[f(w_t)-f(w^*)]
\le
\frac{\|w_0-w^*\|^2}{2c\sqrt{T}}+\frac{c(1+\omega)G^2}{\sqrt{T}} .
\tag{14}
\]

\noindent\textbf{[Step 11]}  By convexity of $f$, the average iterate $\bar w_T$ satisfies
\[
f(\bar w_T)\le\frac{1}{T}\sum_{t=0}^{T-1} f(w_t) .
\]
Taking expectations and applying (14) yields the claimed bound:
\[
\mathbb{E}\bigl[f(\bar w_T)-f(w^*)\bigr]
\le
\frac{\|w_0-w^*\|^2}{2c\sqrt{T}}+\frac{c(1+\omega)G^2}{\sqrt{T}} .
\]
If we replace $(1+\omega)G^2$ by $(L+L\omega+1)G^2$ (using the smoothness inequality to absorb the $L$ term from the $L$‑smoothness constant in a more careful bound) we obtain the statement of Theorem \ref{thm:main}. ∎

\section*{Rate Derivation (With Constants)}
Choosing $c$ that minimizes the right‑hand side of Theorem \ref{thm:main} gives
\[
c^\star = \sqrt{\frac{\|w_0-w^*\|^2}{2(L+L\omega+1)G^2}} .
\]
Plugging $c^\star$ back yields
\[
\mathbb{E}[f(\bar w_T)-f(w^*)]
\le
\sqrt{\frac{2(L+L\omega+1)G^2\|w_0-w^*\|^2}{T}} .
\]
Thus the convergence rate is $\mathcal{O}\!\bigl(T^{-1/2}\bigr)$ with an explicit constant that grows with $\sqrt{1+\omega}$.

\section*{Special Cases}
\begin{itemize}
\item \textbf{Exact gradients, no quantization:} $\omega=0$, $g_t=\nabla f(w_t)$. The bound reduces to the classic SGD rate.
\item \textbf{Deterministic quantization ($\omega=0$ but $Q$ may be biased):} Our proof breaks because unbiasedness (1) is violated; convergence may still hold under a bounded bias assumption but requires a different analysis.
\item \textbf{Strongly convex $f$:} If $f$ is $\mu$‑strongly convex, a constant stepsize $\eta_t=\eta$ with $\eta\le \frac{1}{L(1+\omega)}$ yields a linear convergence term plus a steady‑state error of order $\eta(1+\omega)G^2/\mu$.
\end{itemize}

\end{document}